\section{Asintoto obliquo}
La funzione f ha la retta \(y = mx + q\) come \textbf{asintoto obliquo} per \(x \to +\infty\) se:
\[\lim_{x \to +\infty} (f(x) - mx - q) = 0\]

f ha asintoto obliquo se e solo se:
\[\lim_{x \to +\infty} \frac{f(x)}{x} = m \neq 0 \neq \infty\]
\[\lim_{x \to +\infty} (f(x) - mx) = q \neq \infty\]
\\
Esempio:
\(f(x) = e^x + 2x + 1\) ha asintoto obliquo se \(x \to -\infty\):\\
\[\lim_{x \to -\infty} \frac{e^x + 2x + 1}{x} = 2\]
\[\lim_{x \to -\infty} (e^x + 2x + 1 - 2x) = 1\]

\section{Limite infinito al finito}
\[\lim_{x \to x_0} f(x) = +\infty\]
\[\forall H > 0, \exists \delta > 0 : \forall x \neq x_0, | x - x_0 | < \delta \implies f(x) > H\]
\[\lim_{x \to x_0} f(x) = -\infty\]
\[\forall H < 0, \exists \delta > 0 : \forall x \neq x_0, | x - x_0 | < \delta \implies f(x) < H\]

\section{Asintoto verticale}
\[\lim_{x \to 0^+} \log x = -\infty\]
